\documentclass[12pt,a4paper]{amsart}

\usepackage{amsmath,amssymb,amsthm}
\usepackage{mathtools}
\usepackage{hyperref}
\usepackage{geometry}
\geometry{margin=2.5cm}

% -------------------------------------------------------
\newtheorem{theorem}{Theorem}[section]
\newtheorem{lemma}[theorem]{Lemma}
\newtheorem{corollary}[theorem]{Corollary}
\newtheorem{proposition}[theorem]{Proposition}
\theoremstyle{definition}
\newtheorem{definition}[theorem]{Definition}
\newtheorem{remark}[theorem]{Remark}

\newcommand{\Z}{\mathbb{Z}}
\newcommand{\euler}{\varphi}

% -------------------------------------------------------
\begin{document}
% -------------------------------------------------------

\title{No Semiprime Satisfies Lehmer's Totient Condition:\\ An Arithmetic Proof}

\author{Kurt Ingwer}
\address{Specialist Mathematics Project, Version Build 44}
\email{specialist-maths@localhost}
\date{\today}

\begin{abstract}
Lehmer's totient problem asks whether $\euler(n) \mid (n-1)$ implies $n$ is prime.
We prove that no semiprime $n = p \cdot q$ with $p < q$ prime satisfies this condition.
The key identity is $pq - 1 \equiv p + q - 2 \pmod{(p-1)(q-1)}$: since $(p-1)(q-1) > p + q - 2$
for $p \ge 2$, $q > p$, the quotient $(p-1)(q-1)$ cannot divide $pq - 1$.
This eliminates semiprimes from Lehmer's problem and shows every Lehmer number
(if one exists) has at least three distinct prime factors.
\end{abstract}

\maketitle

% -------------------------------------------------------
\section{Introduction}
% -------------------------------------------------------

Lehmer \cite{Lehmer1932} asked whether $\euler(n) \mid (n-1)$ forces $n$ to be prime.
After establishing that any solution must be squarefree \cite{Ingwer2026sqfree},
the next step is to rule out products of exactly two primes.

For a squarefree $n = p_1 \cdots p_k$, the Lehmer condition becomes
\[
  (p_1-1)(p_2-1)\cdots(p_k-1) \mid p_1 p_2 \cdots p_k - 1.
\]
For $k = 2$ with $p < q$ prime, this is $(p-1)(q-1) \mid pq - 1$.

% -------------------------------------------------------
\section{Main Result}
% -------------------------------------------------------

\begin{theorem}
\label{thm:nosemi}
There is no semiprime $n = p \cdot q$ with $p < q$ prime satisfying
$\euler(n) \mid (n-1)$.
\end{theorem}

\begin{proof}
The condition $\euler(pq) \mid (pq - 1)$ is $(p-1)(q-1) \mid (pq - 1)$.

We use the key identity:
\begin{equation}
  pq - 1 = (p-1)(q-1) + (p-1) + (q-1).
  \label{eq:ident}
\end{equation}

\begin{proof}[Verification of \eqref{eq:ident}]
$(p-1)(q-1) + (p-1) + (q-1) = pq - p - q + 1 + p - 1 + q - 1 = pq - 1$. \checkmark
\end{proof}

From~\eqref{eq:ident}:
\[
  pq - 1 \equiv (p-1) + (q-1) = p + q - 2 \pmod{(p-1)(q-1)}.
\]
For $(p-1)(q-1) \mid (pq-1)$, we need $(p-1)(q-1) \mid (p + q - 2)$.

We show this is impossible:
\begin{equation}
  (p-1)(q-1) > p + q - 2 \quad \text{for } p \ge 2,\ q > p.
  \label{eq:ineq}
\end{equation}

\begin{proof}[Proof of \eqref{eq:ineq}]
We distinguish two cases.

\textbf{Case $p = 2$:} The condition $(p-1)(q-1) \mid (p+q-2)$ becomes $(q-1) \mid q$.
Since $2q - 1 = 2(q-1) + 1$, we get $(q-1) \mid 1$, so $q - 1 = 1$, i.e.\ $q = 2 = p$.
This contradicts $q > p$.

\textbf{Case $p \ge 3$:} We have $p - 2 \ge 1$ and, since $q \ge p + 2 \ge 5$, also $q - 2 \ge 3 \ge 2$.
Note that $(p-1)(q-1) - (p+q-2) = (p-2)(q-2) - 1$.  Therefore
\[
  (p-1)(q-1) - (p+q-2) = (p-2)(q-2) - 1 \ge 1 \cdot 2 - 1 = 1 > 0.
\]
Hence $(p-1)(q-1) > p + q - 2 \ge 1$, and so $(p-1)(q-1) \nmid (p+q-2)$.
\end{proof}

In both cases, $(p-1)(q-1) \nmid (pq-1)$ — a contradiction.
\end{proof}

% -------------------------------------------------------
\section{The Key Identity in Context}
% -------------------------------------------------------

\begin{remark}
Identity~\eqref{eq:ident} generalises: for $n = p_1 \cdots p_k$ squarefree,
\[
  n - 1 = \varphi(n) + \sum_{\emptyset \ne S \subsetneq \{1,\ldots,k\}} \prod_{i \in S}(p_i - 1).
\]
The key modular reduction $n - 1 \equiv \sum_{i=1}^{k}(p_i - 1) \pmod{\euler(n)}$
holds for $k = 2$ and is the basis for the proof above.  For larger~$k$, the
analysis becomes more involved.
\end{remark}

\begin{corollary}
Every Lehmer number has at least three distinct prime factors.
\end{corollary}

\begin{proof}
Squarefreeness \cite{Ingwer2026sqfree} rules out prime powers.
Theorem~\ref{thm:nosemi} rules out products of two distinct primes.
So at least three are needed.
\end{proof}

% -------------------------------------------------------
\begin{thebibliography}{10}

\bibitem{Lehmer1932}
D.~H.~Lehmer,
\emph{On Euler's totient function},
Bull.\ Amer.\ Math.\ Soc.\ \textbf{38} (1932), 745--751.

\bibitem{Cohen1980}
G.~L.~Cohen and P.~Hagis,
\emph{On the number of prime factors of $n$ if $\phi(n) \mid (n-1)$},
Nieuw Arch.\ Wisk.\ (3) \textbf{28} (1980), 177--185.

\bibitem{Ingwer2026sqfree}
K.~Ingwer,
\emph{Every solution to Lehmer's totient problem is squarefree},
Specialist Mathematics Project (2026).

\bibitem{Ingwer2026lehmer3}
K.~Ingwer,
\emph{Lehmer's totient problem for products of three primes},
Specialist Mathematics Project (2026).

\end{thebibliography}

\end{document}
